\chapter{线性泛函}
\label{chap:linear-functionals}
\label{chap:linear-functionals-operators}

研究一个函数时，我们常常只读出一个数：它在某段区间上的平均，它的某个Fourier系数，它在某一点的取值。这样的读数把整个函数变成一个标量。不同读数的可靠程度不同：函数稍有扰动，平均值只变一点，而一个点上的值可能被一个窄尖峰完全改变。哪些读数能由函数的范数稳定地控制，这些读数长什么样，能否从局部扩展到整个空间，是本章的第一组问题。

本章先把函数空间到标量的映射称为泛函，用积分、系数和点值比较线性与非线性读数，再用一个常数界刻画线性泛函的连续性，定义泛函范数。全部连续线性泛函组成对偶空间；在Hilbert空间中它们都是内积，在$L^p$空间中它们都是加权积分。对一般的赋范空间写不出这样的表示，Hahn--Banach定理转而保证子空间上的读数能保持范数地延拓到全空间，由此连续线性读数足以区分不同的向量并读出范数。完备性还给出另一种统一控制：逐个向量有界的泛函族在范数上一致有界。

第二组问题从范数误差转向读数。若一列向量在每个连续线性读数下都收敛，称为弱收敛；在对偶空间中只用原空间的向量去读，称为弱星收敛。我们说明它们与范数收敛的区别、它们仍然控制的范数界，以及在什么条件下有界序列一定有这样收敛的子列。无限维空间中闭球在范数下不紧，第~\ref{chap:functions-as-vectors}章已经看到这一困难；弱星与弱拓扑下的紧性，是绕开它的标准办法。

本章只讨论输出为标量的线性映射。输出为向量的有界线性算子、算子范数、伴随与一般的算子定理在第~\ref{chap:linear-operators}章建立；以泛函为目标的变分与优化由卷五承担。全章的拓扑与度量语言来自第~\ref{chap:topological-spaces}章至第~\ref{chap:metric-space-maps}章，具体空间来自第~\ref{chap:normed-spaces}章和第~\ref{chap:inner-product-spaces}章。

\section{泛函的定义}
\label{sec:lfunc-definition}

函数空间中的元素是整个函数，我们对它所做的许多操作最后只得到一个数。把这类操作统一起来，就得到泛函的概念。本节先看几个常见的例子，逐一核对它们是否线性、在哪个空间上有定义；再研究线性泛函什么时候连续，并用一个数度量它把输入误差放大多少。全章记标量域为$\mathbb F$，它是$\mathbb R$或$\mathbb C$。

\begin{definition}[泛函]
\label{def:lfunc-functional}
设$X$是$\mathbb F$上的线性空间。映射$\ell:X\to\mathbb F$称为$X$上的\term{泛函}（functional）。若对一切$x,y\in X$和$\alpha,\beta\in\mathbb F$都有$\ell(\alpha x+\beta y)=\alpha\ell(x)+\beta\ell(y)$，称$\ell$为\term{线性泛函}（linear functional）。
\end{definition}

有限维空间上的线性泛函由它在一组基上的取值决定，可以写成坐标的线性组合，这在\BookChapterRef{chap:linear-representation}{卷二第2章“线性空间、映射与坐标”}中已经讨论过。函数空间是无限维的，新问题在于：线性泛函未必连续，读数是否稳定要看空间所配的范数。

\subsection{常见泛函}
\label{sec:lfunc-examples}

下面四个例子分别代表加权平均、坐标读取、点上取值和能量。每个例子都要先说清它作用在哪个空间上，因为同一条规则在一个空间上有意义，在另一个空间上可能根本没有定义。

\paragraph{加权积分}
设$(\Omega,\mathcal F,\mu)$是测度空间，$1\leq p\leq\infty$，$q$是\term{共轭指数}（conjugate exponent），即$1/p+1/q=1$，约定$p=1$时$q=\infty$、$p=\infty$时$q=1$。对固定的$g\in L^q(\mu)$，令
\begin{equation}
L_g(f)=\int_\Omega fg\,d\mu\qquad(f\in L^p(\mu)).
\label{eq:lfunc-weighted-integral}
\end{equation}
由第~\ref{chap:normed-spaces}章的Hölder不等式，$\int\lvert fg\rvert\,d\mu\leq\lVert f\rVert_p\lVert g\rVert_q<\infty$，积分有定义。线性来自积分的线性：$L_g(\alpha f_1+\beta f_2)=\alpha L_g(f_1)+\beta L_g(f_2)$。

$L^p$的元素是几乎处处相等的函数组成的等价类，所以还要检查读数不依赖代表的选取。若$f=\tilde f$几乎处处成立，则$fg=\tilde fg$几乎处处成立，两者的积分相同。积分对零测集不敏感，这正是加权积分能在等价类上使用的原因。

\paragraph{正交系数}
设$H$是Hilbert空间，$\{e_k\}$是标准正交系。第$k$个系数$f\mapsto\langle e_k,f\rangle$是$H$上的泛函；按第~\ref{chap:inner-product-spaces}章的约定，内积对第二个变量线性，所以它是线性泛函。在$L^2(0,1)$中，它就是以$\overline{e_k}$为权重的加权积分。

\paragraph{点评估}
设$K$是紧度量空间，$C(K)$是$K$上的连续函数空间。固定$x_0\in K$，令$\delta_{x_0}(f)=f(x_0)$，称为$x_0$处的\term{点评估泛函}（point evaluation functional）。它显然是线性的。

点评估要求逐点取值，只在由真正的函数组成的空间上有定义。在Lebesgue测度下的$L^p(0,1)$中，单点$\{x_0\}$的测度为零，$f$与$f+\mathbf{1}_{\{x_0\}}$属于同一个等价类，在$x_0$处的值却不同。所以“读出$x_0$处的值”不是$L^p$上的映射，谈不上线性或连续。积分与点评估的区别就在这里：前者对零测集不敏感，后者恰恰只看一个零测集。

\paragraph{平方范数}
在赋范空间上，$f\mapsto\lVert f\rVert^2$也把函数变成一个数，例如$L^2$中的能量$\int\lvert f\rvert^2$。它不是线性的：$\lVert2f\rVert^2=4\lVert f\rVert^2\neq2\lVert f\rVert^2$（$f\neq0$时）。非线性泛函在优化中很常见，如以函数为未知量的能量与损失；它们的变分和极小化可参考\BookChapterRef{chap:functional-optimization-variation}{卷五第11章“泛函优化与变分”}。本章只研究线性泛函。

\subsection{泛函范数}
\label{sec:lfunc-norm}

读数是否稳定，就是输入的小误差是否只引起读数的小误差。对线性泛函，这个要求有简单的等价形式。设$X$是赋范空间。由第~\ref{chap:metric-space-maps}章的距离判据，$\ell$在$x_0$处连续，意思是对每个$\varepsilon>0$存在$\delta>0$，使$\lVert x-x_0\rVert<\delta$时$\lvert\ell(x)-\ell(x_0)\rvert<\varepsilon$。

\begin{theorem}[线性泛函的连续性]
\label{thm:lfunc-continuity-bounded}
设$X$是赋范空间，$\ell$是$X$上的线性泛函。以下三条等价：
\begin{enumerate}
\item $\ell$在$0$处连续；
\item $\ell$在$X$上处处连续；
\item 存在常数$C\geq0$，使$\lvert\ell(x)\rvert\leq C\lVert x\rVert$对一切$x\in X$成立。
\end{enumerate}
\end{theorem}

\begin{proof}
由3推2：$\lvert\ell(x)-\ell(x_0)\rvert=\lvert\ell(x-x_0)\rvert\leq C\lVert x-x_0\rVert$，取$\delta=\varepsilon/(C+1)$即可。由2推1是显然的。

由1推3：取$\varepsilon=1$，由$\ell(0)=0$，存在$\delta>0$使$\lVert z\rVert<\delta$时$\lvert\ell(z)\rvert<1$。对$x\neq0$，令$z=\frac{\delta}{2\lVert x\rVert}x$，则$\lVert z\rVert=\delta/2<\delta$，于是由线性
\[
\lvert\ell(x)\rvert=\frac{2\lVert x\rVert}{\delta}\lvert\ell(z)\rvert<\frac{2}{\delta}\lVert x\rVert .
\]
$x=0$时不等式两边都是零。取$C=2/\delta$。
\end{proof}

线性在证明中起了两次作用：把任意点处的误差平移成原点附近的误差，把原点附近的控制按比例放大到所有向量。满足第3条的线性泛函称为\term{有界线性泛函}（bounded linear functional）；由定理~\ref{thm:lfunc-continuity-bounded}，对线性泛函而言“有界”与“连续”是同一回事。这里的有界指在单位球上有界，而不是在整个空间上有界，后者只有零泛函才满足。

\begin{definition}[泛函范数]
\label{def:lfunc-functional-norm}
设$\ell$是赋范空间$X$上的有界线性泛函，令
\[
\lVert\ell\rVert=\sup_{\lVert x\rVert\leq1}\lvert\ell(x)\rvert,
\]
称为$\ell$的\term{泛函范数}（norm of a functional）。
\end{definition}

由线性，$\lvert\ell(x)\rvert\leq\lVert\ell\rVert\lVert x\rVert$对一切$x$成立，并且$\lVert\ell\rVert$是使这一不等式成立的最小常数$C$。它衡量的是读数对输入误差的最大放大倍数。

\paragraph{点评估的范数依赖空间}
在$C([0,1])$上配一致范数$\lVert f\rVert_\infty=\max\lvert f\rvert$，由$\lvert f(0)\rvert\leq\lVert f\rVert_\infty$得$\lVert\delta_0\rVert\leq1$，取$f\equiv1$得$\lVert\delta_0\rVert=1$。点评估在一致范数下是稳定的读数。

同一个集合$C([0,1])$改配积分范数，结论完全不同。取贴着左端点的尖峰
\[
\tau_n(t)=\max(1-nt,0)\qquad(t\in[0,1]),
\]
它连续，$\tau_n(0)=1$，而直接积分得
\[
\lVert\tau_n\rVert_1=\int_0^{1/n}(1-nt)\,dt=\frac{1}{2n},\qquad
\lVert\tau_n\rVert_2^2=\int_0^{1/n}(1-nt)^2\,dt=\frac{1}{3n}.
\]
两种积分范数都趋于零，读数却始终是$1$，所以不存在常数$C$使$\lvert f(0)\rvert\leq C\lVert f\rVert_2$对一切连续$f$成立。在$(C([0,1]),\lVert\cdot\rVert_2)$上，$\delta_0$是线性的、处处有定义，却不连续。

\paragraph{加权积分的范数也依赖空间}
在$L^2(0,1)$上取$g=\mathbf{1}_{[0,1/2]}$，$L_g(f)=\int_0^{1/2}f$。Hölder不等式给出$\lVert L_g\rVert\leq\lVert g\rVert_2=1/\sqrt2$；取$f=g$，$L_g(g)=1/2$而$\lVert g\rVert_2=1/\sqrt2$，比值为$1/\sqrt2$，所以$\lVert L_g\rVert=1/\sqrt2$。同一规则在$C([0,1])$的一致范数下，由$\lvert\int_0^{1/2}f\rvert\leq\lVert f\rVert_\infty/2$及$f\equiv1$得范数为$1/2$。读数规则相同，范数不同，因为衡量输入误差的方式不同。是否稳定、放大多少，都依赖所选的空间与范数。

\section{泛函的表示}
\label{sec:lfunc-representation}

上一节逐个研究读数规则。现在把同一空间上全部连续线性读数收集起来，看它们组成怎样的空间，以及在具体空间中它们能否写成熟悉的形式。先建立对偶空间及其范数，并证明它总是完备的；再分别在Hilbert空间和$L^p$空间中证明，每个连续线性泛函都是内积或加权积分，表示唯一，并且泛函范数等于表示对象的范数。

\subsection{对偶空间}
\label{sec:lfunc-dual-space}

两个有界线性泛函之和、数乘仍是有界线性泛函，所以它们组成一个线性空间，泛函范数是这个空间上的范数：三角不等式来自$\lvert(\ell_1+\ell_2)(x)\rvert\leq\lvert\ell_1(x)\rvert+\lvert\ell_2(x)\rvert$再对单位球取上确界。

\begin{definition}[对偶空间与二次对偶空间]
\label{def:lfunc-dual-space}
赋范空间$X$上全部有界线性泛函组成的空间记为$X^*$，配以泛函范数，称为$X$的\term{对偶空间}（dual space）。$X^*$本身是赋范空间，它的对偶空间$X^{**}=(X^*)^*$称为$X$的\term{二次对偶空间}（bidual space）。
\end{definition}

对偶空间的元素是读数规则，二次对偶空间的元素是“读数规则的读数规则”：给它一个$X$上的泛函，它返回一个数。第~\ref{sec:lfunc-reflexivity}小节会看到，每个$x\in X$都以$\ell\mapsto\ell(x)$的方式自然地成为$X^{**}$的元素；比较这些元素与$X^{**}$的全部元素，是后面讨论弱拓扑与紧性的基础。

\paragraph{对偶范数可以不同于原范数}
即使在有限维中，对偶空间的范数也可以与原空间的范数形式不同。在$\mathbb R^2$上配$\lVert\bm x\rVert_1=\lvert x_1\rvert+\lvert x_2\rvert$。每个线性泛函都写成$\ell(\bm x)=a_1x_1+a_2x_2$，由
\[
\lvert\ell(\bm x)\rvert\leq\max(\lvert a_1\rvert,\lvert a_2\rvert)\bigl(\lvert x_1\rvert+\lvert x_2\rvert\bigr)
\]
且在$\bm x=\bm e_1$或$\bm e_2$处取到等号，得$\lVert\ell\rVert=\max(\lvert a_1\rvert,\lvert a_2\rvert)$。系数向量$\bm a$作为读数的长度要用最大值范数计算，这是\term{对偶范数}（dual norm）的一个例子。所以不能把读数规则与表示它的系数向量随意等同，除非已经核对了范数。

\paragraph{对偶空间总是完备的}
对偶空间有一个原空间未必具有的性质：无论$X$本身是否完备，$X^*$都完备。原因是读数取值在完备的$\mathbb F$中。

\begin{theorem}[对偶空间的完备性]
\label{thm:lfunc-dual-complete}
对任意赋范空间$X$，$X^*$是Banach空间。
\end{theorem}

\begin{proof}
设$(\ell_n)$是$X^*$中的Cauchy列。对每个$x$，$\lvert\ell_n(x)-\ell_m(x)\rvert\leq\lVert\ell_n-\ell_m\rVert\lVert x\rVert$，所以$(\ell_n(x))$是$\mathbb F$中的Cauchy列，有极限，记为$\ell(x)$。极限保持线性，$\ell$是线性泛函。

给定$\varepsilon>0$，取$N$使$n,m\geq N$时$\lVert\ell_n-\ell_m\rVert<\varepsilon$。对$\lVert x\rVert\leq1$，$\lvert\ell_n(x)-\ell_m(x)\rvert<\varepsilon$；令$m\to\infty$得$\lvert\ell_n(x)-\ell(x)\rvert\leq\varepsilon$。所以$\lvert\ell(x)\rvert\leq\lVert\ell_N\rVert+\varepsilon$在单位球上成立，$\ell$有界，属于$X^*$；并且$n\geq N$时$\lVert\ell_n-\ell\rVert\leq\varepsilon$，即$\ell_n\to\ell$。
\end{proof}

例如由多项式组成的不完备空间配上一致范数，它的对偶空间仍是Banach空间。这一性质在第~\ref{sec:lfunc-uniform-boundedness}节和第~\ref{sec:lfunc-weak-norm-control}小节中起关键作用：需要完备性的论证，可以在$X^*$上进行而不必要求$X$完备。

\subsection{Hilbert空间的泛函表示}
\label{sec:lfunc-riesz}

在Hilbert空间$H$中，每个向量都给出一个读数。按第~\ref{chap:inner-product-spaces}章的约定，复内积对第一个变量共轭线性、对第二个变量线性。对$h\in H$令$\ell_h(f)=\langle h,f\rangle$，它对$f$线性；由Cauchy--Schwarz不等式，$\lvert\ell_h(f)\rvert\leq\lVert h\rVert\lVert f\rVert$，所以$\lVert\ell_h\rVert\leq\lVert h\rVert$。$h\neq0$时取$f=h/\lVert h\rVert$，读数为$\lVert h\rVert$，于是$\lVert\ell_h\rVert=\lVert h\rVert$。

反过来的问题更有用：给定连续线性泛函$\ell$，是否一定有向量$h$使$\ell=\ell_h$？若是，读数就可以用空间中的几何来研究，读数的范数就是一个向量的长度。

\paragraph{零空间与正交方向}
零泛函由零向量表示。设$\ell\neq0$，考察它的零空间$N=\{f:\ell(f)=0\}$。它是线性子空间；又由连续性，若$f_n\in N$且$f_n\to f$，则$\ell(f)=\lim\ell(f_n)=0$，所以$N$是闭子空间，且$N\neq H$。若能找到与$N$正交的非零向量$w$，读数就只取决于$f$沿$w$方向的分量，表示向量应当是$w$的某个倍数。由此得到的结论称为\term{Riesz表示定理}（Riesz representation theorem）。

\begin{theorem}[Riesz表示定理]
\label{thm:lfunc-riesz}
设$H$是Hilbert空间，$\ell\in H^*$。则存在唯一的$h_\ell\in H$，使
\begin{equation}
\ell(f)=\langle h_\ell,f\rangle\qquad(f\in H),
\label{eq:lfunc-riesz}
\end{equation}
并且$\lVert\ell\rVert=\lVert h_\ell\rVert$。
\end{theorem}

\begin{proof}
$\ell=0$时取$h_\ell=0$。下设$\ell\neq0$，$N$如上。

\emph{找与零空间正交的方向。}取$v$使$\ell(v)\neq0$。由第~\ref{chap:inner-product-spaces}章的正交投影定理，$v=Pv+w$，其中$Pv\in N$，$w\perp N$。由$\ell(w)=\ell(v)-\ell(Pv)=\ell(v)\neq0$知$w\neq0$。

\emph{写出表示。}对任意$f\in H$，向量$f-\frac{\ell(f)}{\ell(w)}w$在$\ell$下的值为零，属于$N$，因而与$w$正交：
\[
0=\Bigl\langle w,f-\frac{\ell(f)}{\ell(w)}w\Bigr\rangle=\langle w,f\rangle-\frac{\ell(f)}{\ell(w)}\lVert w\rVert^2 .
\]
这里用到内积对第二个变量线性。解出$\ell(f)=\frac{\ell(w)}{\lVert w\rVert^2}\langle w,f\rangle$。令$h_\ell=\frac{\overline{\ell(w)}}{\lVert w\rVert^2}w$，由内积对第一个变量共轭线性，$\langle h_\ell,f\rangle=\frac{\ell(w)}{\lVert w\rVert^2}\langle w,f\rangle=\ell(f)$。

\emph{唯一性与范数。}若$h,h'$都满足式~\eqref{eq:lfunc-riesz}，则$\langle h-h',f\rangle=0$对一切$f$成立，取$f=h-h'$得$h=h'$。范数等式是本小节开头对$\ell_h$的计算。
\end{proof}

证明只用了正交投影定理，而正交投影定理对任意Hilbert空间中的闭子空间成立，所以Riesz表示定理不要求$H$可分，也不需要事先选取正交基。表示向量公式中的共轭使对应$\ell\mapsto h_\ell$在复空间中是共轭线性的：$h_{\alpha\ell}=\bar\alpha h_\ell$。实空间中这个共轭消失，$H^*$与$H$通过保范数的线性双射等同。

\paragraph{\texorpdfstring{$L^2$}{L2}中的读数与方向}
在$L^2(0,1)$中，几个读数的表示向量可以直接写出。整体平均$\ell(f)=\int_0^1f$对应$h_\ell=1$，范数为$1$。加权读数$\ell(f)=\int_0^1tf(t)\,dt$对应$h_\ell(t)=t$，范数为$(\int_0^1t^2\,dt)^{1/2}=1/\sqrt3$。第$k$个正交系数$f\mapsto\langle e_k,f\rangle$对应$h_\ell=e_k$。

\paragraph{表示以连续性为前提}
定理不能替一条读数规则证明连续性。点评估在$L^2(0,1)$上根本没有定义；在$(C([0,1]),\lVert\cdot\rVert_2)$上它有定义但不连续，也确实没有表示向量：若有$h$使$f(0)=\langle h,f\rangle$对一切连续$f$成立，代入尖峰$\tau_n$，左端恒为$1$，右端不超过$\lVert h\rVert_2\lVert\tau_n\rVert_2\to0$，矛盾。完备性也不能省：$(C([0,1]),\lVert\cdot\rVert_2)$上的有界读数$\int_0^{1/2}f$在完备化$L^2(0,1)$中的表示向量是$\mathbf{1}_{[0,1/2]}$，它没有连续代表，所以在不完备的空间中找不到表示向量。

若某个函数空间配上内积后完备，并且点评估在其中连续，Riesz表示定理就把点评估写成与某个固定函数的内积。第~\ref{chap:rkhs}章研究的再生核Hilbert空间正是这样的空间。

\subsection{\texorpdfstring{$L^p$}{Lp}空间的泛函表示}
\label{sec:lfunc-lp-dual}

在$L^p$中没有内积可用，自然的候选是加权积分~\eqref{eq:lfunc-weighted-integral}。第~\ref{sec:lfunc-examples}小节已说明每个$g\in L^q$给出有界读数$L_g$，且$\lVert L_g\rVert\leq\lVert g\rVert_q$。本小节回答反方向的问题。为了不处理共轭记号，下面在实值函数上证明；复值情形对实部与虚部分别处理，结论相同，只是$p=2$时$L_g$的表示向量在内积记号下是$\bar g$。

\begin{theorem}[$L^p$的对偶表示]
\label{thm:lfunc-lp-dual}
设$(\Omega,\mathcal F,\mu)$是$\sigma$有限测度空间，$1\leq p<\infty$，$q$为共轭指数。则每个$\ell\in(L^p(\mu))^*$都存在$g\in L^q(\mu)$使$\ell=L_g$；$g$在几乎处处相等的意义下唯一，并且$\lVert\ell\rVert=\lVert g\rVert_q$。
\end{theorem}

证明分三步：先说明$\lVert L_g\rVert$恰好等于$\lVert g\rVert_q$，从而表示唯一；再在有限测度下由读数构造测度并用Radon--Nikodym定理得到$g$；最后推广到$\sigma$有限测度。

\paragraph{范数等式与唯一性}
设$1<p<\infty$，$g\neq0$。令$f=\lvert g\rvert^{q-1}\operatorname{sgn}g$，则$fg=\lvert g\rvert^q$，又因$(q-1)p=q$，$\lvert f\rvert^p=\lvert g\rvert^q$。于是$L_g(f)=\lVert g\rVert_q^q$，$\lVert f\rVert_p=\lVert g\rVert_q^{q/p}$，比值为$\lVert g\rVert_q^{q-q/p}=\lVert g\rVert_q$。$p=1$、$q=\infty$时等号未必取到，改用逼近：对$0<\varepsilon<\lVert g\rVert_\infty$，集合$E=\{\lvert g\rvert>\lVert g\rVert_\infty-\varepsilon\}$测度为正；由$\sigma$有限，可取$E$中测度有限且为正的部分$E'$，令$f=\operatorname{sgn}(g)\mathbf{1}_{E'}/\mu(E')$，则$\lVert f\rVert_1=1$而$L_g(f)\geq\lVert g\rVert_\infty-\varepsilon$。两种情形都得到$\lVert L_g\rVert=\lVert g\rVert_q$。

由此表示是唯一的：若$L_g=L_{g'}$，则$L_{g-g'}=0$，其范数$\lVert g-g'\rVert_q=0$，即$g=g'$几乎处处成立。表示函数只确定到几乎处处相等，这与$L^q$的元素本来就是等价类一致。

\paragraph{有限测度下构造表示函数}
先设$\mu(\Omega)<\infty$。给定$\ell\in(L^p)^*$，对可测集$E$令$\nu(E)=\ell(\mathbf{1}_E)$。有限可加性来自线性。可数可加性来自连续性，这里用到$p<\infty$：若$E=\bigcup_kE_k$为不交并，则
\[
\Bigl\lVert\mathbf{1}_E-\sum_{k\leq n}\mathbf{1}_{E_k}\Bigr\rVert_p=\mu\Bigl(\bigcup_{k>n}E_k\Bigr)^{1/p}\to0,
\]
所以$\sum_{k\leq n}\nu(E_k)\to\nu(E)$。又$\lvert\nu(E)\rvert\leq\lVert\ell\rVert\mu(E)^{1/p}$，$\nu$是有限的有符号测度，并且$\mu(E)=0$时$\nu(E)=0$，即$\nu$对$\mu$绝对连续。

由第~\ref{chap:measure}章，有限的有符号测度可分解为两个有限测度之差，它们都对$\mu$绝对连续；分别应用第~\ref{chap:measure-densities}章的Radon--Nikodym定理，得到$g\in L^1(\mu)$使$\ell(\mathbf{1}_E)=\int_Eg\,d\mu$对一切可测集$E$成立。由线性，$\ell(f)=\int fg\,d\mu$对一切简单函数$f$成立。

\paragraph{表示函数属于\texorpdfstring{$L^q$}{Lq}}
设$1<p<\infty$。令$E_n=\{\lvert g\rvert\leq n\}$，$f_n=\lvert g\rvert^{q-1}\operatorname{sgn}(g)\mathbf{1}_{E_n}$。$f_n$是有界可测函数，可以被简单函数一致逼近；在有限测度下一致收敛推出$L^p$收敛，而$\ell$连续、$g$可积，所以等式$\ell(f_n)=\int f_ng\,d\mu$也成立。于是
\[
\int_{E_n}\lvert g\rvert^q\,d\mu=\ell(f_n)\leq\lVert\ell\rVert\lVert f_n\rVert_p=\lVert\ell\rVert\Bigl(\int_{E_n}\lvert g\rvert^q\,d\mu\Bigr)^{1/p},
\]
整理得$(\int_{E_n}\lvert g\rvert^q)^{1/q}\leq\lVert\ell\rVert$，由单调收敛定理令$n\to\infty$得$\lVert g\rVert_q\leq\lVert\ell\rVert$。

设$p=1$。若$F=\{g>\lVert\ell\rVert\}$测度为正，取$f=\mathbf{1}_F/\mu(F)$，则$\lVert f\rVert_1=1$，而$\ell(f)=\frac{1}{\mu(F)}\int_Fg\,d\mu>\lVert\ell\rVert$，矛盾；对$\{g<-\lVert\ell\rVert\}$同理。所以$\lVert g\rVert_\infty\leq\lVert\ell\rVert$。

两种情形下$g\in L^q$，$\ell$与$L_g$都是$L^p$上的连续线性泛函，在简单函数上相等。由第~\ref{chap:normed-spaces}章，简单函数在$L^p$（$p<\infty$）中稠密，两个连续函数在稠密集上相等则处处相等，所以$\ell=L_g$。

\paragraph{推广到\texorpdfstring{$\sigma$}{sigma}有限测度}
设$\Omega=\bigcup_k\Omega_k$，$\Omega_k$两两不交且测度有限。把$\ell$限制到在$\Omega_k$外为零的函数上，由有限测度情形得到$\Omega_k$上的$g_k$，令$g=\sum_kg_k\mathbf{1}_{\Omega_k}$。对$A_n=\Omega_1\cup\cdots\cup\Omega_n$，同样的论证给出$\lVert g\mathbf{1}_{A_n}\rVert_q\leq\lVert\ell\rVert$，由单调收敛得$\lVert g\rVert_q\leq\lVert\ell\rVert$。于是$\ell$与$L_g$在支撑于某个$A_n$内的简单函数上相等，这些函数在$L^p$中仍稠密，所以$\ell=L_g$。$p=1$时$\sigma$有限的假设不能去掉；$1<p<\infty$时可以用更细的论证去掉，这里不展开。

\begin{example}[一个可复算的对偶范数]
\label{ex:lfunc-l3-norm}
在$L^3(0,1)$上取$\ell(f)=\int_0^1tf(t)\,dt$。共轭指数$q=3/2$，由定理~\ref{thm:lfunc-lp-dual}，
\[
\lVert\ell\rVert=\Bigl(\int_0^1t^{3/2}\,dt\Bigr)^{2/3}=\Bigl(\frac25\Bigr)^{2/3}\approx0.543,
\]
取到范数的输入是$f(t)=t^{q-1}=t^{1/2}$：$\ell(t^{1/2})=\int_0^1t^{3/2}=2/5$，$\lVert t^{1/2}\rVert_3=(\int_0^1t^{3/2})^{1/3}=(2/5)^{1/3}$，比值为$(2/5)^{2/3}$。同一个权重在$L^2$中的读数范数是$\lVert t\rVert_2=1/\sqrt3\approx0.577$，在$L^1$中是$\lVert t\rVert_\infty=1$。
\end{example}

\paragraph{端点\texorpdfstring{$p=\infty$}{p=infinity}}
定理排除了$p=\infty$，这不是技术上的疏忽。每个$g\in L^1$仍给出$L^\infty$上的有界读数$L_g$，但$L^\infty$上的连续线性泛函一般不能全部写成这种形式。第~\ref{sec:lfunc-hahn-banach}小节将用延拓定理在$\ell^\infty$上构造一个不来自任何$\ell^1$序列的连续线性泛函。所以有限维和$L^2$中“对偶就是自身”的直觉，以及$1\leq p<\infty$时“对偶就是$L^q$”的结论，都不能直接用到$L^\infty$上。

\section{泛函的延拓}
\label{sec:lfunc-extension}

一般的赋范空间没有内积，也不是积分空间，上一节的表示定理都用不上，我们甚至不知道$X^*$中是否有足够多的元素。在Hilbert空间中，任何非零向量$x$都被读数$\langle x,\cdot\rangle$识别出来，因为它在$x$处的值是$\lVert x\rVert^2\neq0$。一般空间中写不出这样的读数，但可以先在一个小子空间上定义它，再设法扩展到整个空间，并且不增大它的范数。本节先建立这种延拓，再用它证明连续线性读数能读出范数、区分向量。

\subsection{Hahn--Banach延拓}
\label{sec:lfunc-hahn-banach}

设$M$是$X$的线性子空间，$\ell$是$M$上的有界线性泛函，即$\lvert\ell(m)\rvert\leq\lVert\ell\rVert\lVert m\rVert$对$m\in M$成立。我们要找$L\in X^*$，使$L$在$M$上与$\ell$一致，并且$\lVert L\rVert=\lVert\ell\rVert$。延拓总能保持在$M$上的取值，难点在于不增大范数：$L$在$M$外的取值要受$\lvert L(x)\rvert\leq\lVert\ell\rVert\lVert x\rVert$的约束。

\begin{theorem}[Hahn--Banach定理]
\label{thm:lfunc-hahn-banach}
设$X$是赋范空间，$M\subset X$是线性子空间，$\ell:M\to\mathbb F$是有界线性泛函。则存在$L\in X^*$，使$L|_M=\ell$且$\lVert L\rVert=\lVert\ell\rVert$。
\end{theorem}

这一结论称为\term{Hahn--Banach定理}（Hahn--Banach theorem）。下面先完整证明延拓一维的核心步骤，再说明怎样推广到整个空间。

\paragraph{延拓一维}
设$\mathbb F=\mathbb R$，不妨设$\lVert\ell\rVert=1$（$\ell=0$时取$L=0$，否则除以$\lVert\ell\rVert$）。任取$z\notin M$，要把$\ell$延拓到$M\oplus\operatorname{span}\{z\}$上。线性延拓完全由一个数$t=L(z)$决定：$L(m+sz)=\ell(m)+st$。要求$\lVert L\rVert\leq1$，即
\[
\lvert\ell(m)+st\rvert\leq\lVert m+sz\rVert\qquad(m\in M,\ s\in\mathbb R).
\]
$s=0$时这就是$\lVert\ell\rVert\leq1$。$s\neq0$时两边除以$\lvert s\rvert$，并把$m/s$重新记为$m$，条件化为
\[
-\lVert m+z\rVert-\ell(m)\leq t\leq\lVert m+z\rVert-\ell(m)\qquad(m\in M).
\]

这样的$t$存在，当且仅当左端对$m$的上确界不超过右端对$m'$的下确界，即对一切$m,m'\in M$，
\[
-\lVert m+z\rVert-\ell(m)\leq\lVert m'+z\rVert-\ell(m').
\]
这确实成立：$\ell(m')-\ell(m)=\ell(m'-m)\leq\lVert m'-m\rVert\leq\lVert m'+z\rVert+\lVert m+z\rVert$，移项即得。所以可选的$t$组成一个非空闭区间，任取其中一个，就把$\ell$保范数地延拓了一维。

\paragraph{从一维到全空间}
若$X$可分，取可数稠密集$\{z_1,z_2,\ldots\}$，逐个把$z_k$添入定义域并重复上一步，得到$\operatorname{span}(M\cup\{z_k\}_k)$上范数仍为$\lVert\ell\rVert$的延拓。这个子空间稠密，而有界线性泛函一致连续，可以唯一地连续延拓到闭包即整个$X$上，范数不变：若$x_n\to x$，则$(L(x_n))$是Cauchy列，极限与所选序列无关，且$\lvert L(x)\rvert=\lim\lvert L(x_n)\rvert\leq\lVert\ell\rVert\lVert x\rVert$。不可分时，“逐个添加”要换成对所有保范数延拓按包含关系排序后使用Zorn引理，思路相同，这里不给出这一部分的证明。

\paragraph{复空间}
设$\mathbb F=\mathbb C$。复线性泛函由实部决定：记$u=\operatorname{Re}\ell$，则$\ell(m)=u(m)-\mathrm iu(\mathrm im)$，因为$\operatorname{Im}\ell(m)=-\operatorname{Re}(\mathrm i\ell(m))=-u(\mathrm im)$。把$X$看成实赋范空间，$u$是$M$上的实线性泛函，$\lvert u(m)\rvert\leq\lvert\ell(m)\rvert$，按实情形把$u$延拓为$U$且$\lVert U\rVert\leq\lVert\ell\rVert$，再令$L(x)=U(x)-\mathrm iU(\mathrm ix)$。直接验证$L$复线性且在$M$上等于$\ell$。对每个$x$取$\theta$使$e^{-\mathrm i\theta}L(x)=\lvert L(x)\rvert$，则$\lvert L(x)\rvert=L(e^{-\mathrm i\theta}x)$是实数，等于其实部$U(e^{-\mathrm i\theta}x)\leq\lVert\ell\rVert\lVert x\rVert$。所以$\lVert L\rVert\leq\lVert\ell\rVert$，反向不等式是显然的。

\paragraph{延拓可以不唯一}
一维延拓中的“非空区间”一般不止一点。在$(\mathbb R^2,\lVert\cdot\rVert_1)$中取$M$为横轴，$\ell(s,0)=s$，$\lVert\ell\rVert=1$。延拓写成$L_t(\bm x)=x_1+tx_2$，由第~\ref{sec:lfunc-dual-space}小节的对偶范数计算，$\lVert L_t\rVert=\max(1,\lvert t\rvert)$，所以每个$\lvert t\rvert\leq1$都给出保范数延拓。几何上，直线$x_1+tx_2=1$都经过$(1,0)$并支撑菱形单位球，菱形在这个顶点处有许多支撑线。换成欧氏范数，单位圆在$(1,0)$只有一条切线$x_1=1$，延拓唯一。定理保证的是存在，并且保持的只有两条约束：在$M$上的取值和范数上界，$M$外的取值一般不由它们确定。

\paragraph{区分向量与子空间}
延拓的一个直接用途，是用一个连续线性读数把一个向量与一个闭子空间分开。

\begin{theorem}[距离泛函]
\label{thm:lfunc-distance-functional}
设$M$是赋范空间$X$的闭子空间，$x_0\notin M$，$d=\inf_{m\in M}\lVert x_0-m\rVert$。则$d>0$，并且存在$L\in X^*$，使$\lVert L\rVert=1$，$L$在$M$上为零，$L(x_0)=d$。
\end{theorem}

\begin{proof}
$M$闭而$x_0\notin M$，所以$x_0$有一个邻域与$M$不交，$d>0$。在$M\oplus\operatorname{span}\{x_0\}$上令$\ell(m+sx_0)=sd$。$s\neq0$时$\lVert m+sx_0\rVert=\lvert s\rvert\lVert x_0+m/s\rVert\geq\lvert s\rvert d$，所以$\lvert\ell(m+sx_0)\rvert\leq\lVert m+sx_0\rVert$，$\lVert\ell\rVert\leq1$。取$m_n\in M$使$\lVert x_0-m_n\rVert\to d$，则$\ell(x_0-m_n)=d$，比值趋于$1$，故$\lVert\ell\rVert=1$。由定理~\ref{thm:lfunc-hahn-banach}延拓即得。
\end{proof}

\paragraph{\texorpdfstring{$\ell^\infty$}{l-infinity}上不来自\texorpdfstring{$\ell^1$}{l1}的泛函}
现在可以完成第~\ref{sec:lfunc-lp-dual}小节末尾的说明。记$c$为$\ell^\infty$中收敛数列组成的子空间，在其上令$\lambda(x)=\lim_nx_n$。由$\lvert\lim x_n\rvert\leq\sup\lvert x_n\rvert$及常数列，$\lVert\lambda\rVert=1$。由定理~\ref{thm:lfunc-hahn-banach}，$\lambda$延拓为某个$\Lambda\in(\ell^\infty)^*$。

若$\Lambda$来自某个$y\in\ell^1$，即$\Lambda(x)=\sum_ky_kx_k$，则对标准单位向量$\bm e_k$（只有第$k$项为$1$的数列）有$y_k=\Lambda(\bm e_k)=\lim_n(\bm e_k)_n=0$，于是$y=0$、$\Lambda=0$；可是$\Lambda$在常数列$(1,1,\ldots)$上的值为$1$，矛盾。所以$(\ell^\infty)^*$严格大于由$\ell^1$给出的那部分。这个$\Lambda$没有显式公式，它的存在完全依靠延拓定理；后面讨论弱收敛与反射性时还会再用到它。

\subsection{范数的对偶表示}
\label{sec:lfunc-norm-dual}

Hilbert空间中，$x\neq0$时单位读数$\langle x/\lVert x\rVert,\cdot\rangle$在$x$处恰好读出$\lVert x\rVert$。在一般赋范空间中，我们用延拓构造一个起同样作用的读数：先在$x$张成的一维子空间上规定它读出$\lVert x\rVert$，再保范数地延拓。

\begin{theorem}[范数的对偶表示]
\label{thm:lfunc-norming}
设$X$是赋范空间，$x\in X$，$x\neq0$。则存在$L\in X^*$，使$\lVert L\rVert=1$且$L(x)=\lVert x\rVert$。因此对一切$x\in X$，
\begin{equation}
\lVert x\rVert=\max_{\lVert L\rVert\leq1}\lvert L(x)\rvert ,
\label{eq:lfunc-norm-dual}
\end{equation}
并且若$L(x)=L(y)$对一切$L\in X^*$成立，则$x=y$。
\end{theorem}

\begin{proof}
在$M=\operatorname{span}\{x\}$上令$\ell(\alpha x)=\alpha\lVert x\rVert$，则$\lvert\ell(\alpha x)\rvert=\lVert\alpha x\rVert$，$\lVert\ell\rVert=1$。由定理~\ref{thm:lfunc-hahn-banach}延拓得到所需的$L$。另一方面，对一切$\lVert L'\rVert\leq1$有$\lvert L'(x)\rvert\leq\lVert x\rVert$，而上面的$L$取到等号，所以式~\eqref{eq:lfunc-norm-dual}中的上确界可以取到，是最大值；$x=0$时两边为零。最后，若$x\neq y$，对$x-y$取这样的$L$，得$L(x)-L(y)=\lVert x-y\rVert\neq0$。
\end{proof}

式~\eqref{eq:lfunc-norm-dual}说，一个向量的长度可以由全部单位读数的最大值恢复，向量本身也由全部读数唯一确定。这里的范数与对偶范数是对称的：泛函范数是单位球上读数的上确界，向量范数是单位泛函读数的最大值。

\paragraph{具体空间中的取范读数}
在具体空间中，取到最大值的读数常常可以直接写出。在$C([0,1])$中，若$\lvert f\rvert$在$t_0$处取到最大值，则$L=\overline{\operatorname{sgn}f(t_0)}\,\delta_{t_0}$满足$\lVert L\rVert=1$、$L(f)=\lVert f\rVert_\infty$，这里$\operatorname{sgn}z=z/\lvert z\rvert$。在实值$L^p$（$1<p<\infty$）中，对$f\neq0$取权重$g=\lvert f\rvert^{p-1}\operatorname{sgn}f/\lVert f\rVert_p^{p-1}$，第~\ref{sec:lfunc-lp-dual}小节的计算给出$\lVert g\rVert_q=1$、$L_g(f)=\lVert f\rVert_p$。定理的价值在于，即使写不出这样的读数，它也保证存在。

凸集之间的分离、优化问题的对偶和最优性条件，都是同一延拓思想的几何形式，属于优化问题，可参考\BookChapterRef{chap:constrained-optimization}{卷五第5章“约束优化与对偶”}。

\section{泛函族的一致有界性}
\label{sec:lfunc-uniform-boundedness}

设$\{\ell_\alpha\}_{\alpha\in I}$是$X^*$中的一族泛函，并且在每个固定的向量上读数都有界：对每个$x$，$\sup_\alpha\lvert\ell_\alpha(x)\rvert<\infty$。这称为逐点有界。我们想知道能否推出$\sup_\alpha\lVert\ell_\alpha\rVert<\infty$，即一致有界。两者的量词次序不同：前者是“对每个$x$，存在常数$C_x$”，后者是“存在常数$C$，对每个单位向量$x$”。下一节研究弱收敛时，正需要从逐个读数的信息得到这样的统一控制。

\paragraph{不完备时可以失败}
在只有有限项非零的数列组成的空间$c_{00}$上配最大值范数，令$\ell_n(x)=nx_n$。对固定的$x$，$n$超过它最后一个非零坐标后$\ell_n(x)=0$，所以$\sup_n\lvert\ell_n(x)\rvert<\infty$；可是在$\bm e_n$上读数为$n$，$\lVert\ell_n\rVert=n$，没有统一上界。$c_{00}$不完备，例如$(1,1/2,\ldots,1/n,0,0,\ldots)$是Cauchy列，极限$(1/n)_n$不在$c_{00}$中。下面会看到，完备性正是把逐点有界转为一致有界的关键。

\paragraph{需要的工具}
第~\ref{chap:metric-spaces}章的Baire纲定理说：在非空完备度量空间中，可数个内部为空的闭集不能覆盖整个空间。换句话说，若非空完备度量空间写成可数个闭集之并，则至少有一个闭集含有某个开球。证明的思路是用嵌套闭球逐个避开各个闭集，完备性保证球心的极限存在；我们直接使用这一结论。

\begin{theorem}[泛函的一致有界原理]
\label{thm:lfunc-uniform-boundedness}
设$X$是Banach空间，$\{\ell_\alpha\}_{\alpha\in I}\subset X^*$，并且对每个$x\in X$有$\sup_\alpha\lvert\ell_\alpha(x)\rvert<\infty$。则$\sup_\alpha\lVert\ell_\alpha\rVert<\infty$。
\end{theorem}

这一结论称为\term{一致有界原理}（uniform boundedness principle）。证明按“闭集覆盖、找内点、平移到原点、缩放到单位球”的顺序进行。

\begin{proof}
\emph{闭集覆盖。}对正整数$k$令$F_k=\{x\in X:\sup_\alpha\lvert\ell_\alpha(x)\rvert\leq k\}=\bigcap_\alpha\{x:\lvert\ell_\alpha(x)\rvert\leq k\}$。每个$\ell_\alpha$连续，所以交集中的每个集合都是闭集，$F_k$也是闭集。逐点有界说明每个$x$属于某个$F_k$，于是$X=\bigcup_kF_k$。

\emph{找内点。}$X$完备且非空，由Baire纲定理，存在$k_0$和开球$B(x_0,r)\subset F_{k_0}$。

\emph{平移到原点。}对任意$\lVert z\rVert<r$，$x_0+z$与$x_0$都在$F_{k_0}$中，所以对每个$\alpha$，
\[
\lvert\ell_\alpha(z)\rvert=\lvert\ell_\alpha(x_0+z)-\ell_\alpha(x_0)\rvert\leq2k_0 .
\]

\emph{缩放到单位球。}对单位向量$x$及$0<s<r$，$\lVert sx\rVert<r$，由线性$\lvert\ell_\alpha(x)\rvert=s^{-1}\lvert\ell_\alpha(sx)\rvert\leq2k_0/s$。令$s\to r$，得$\lVert\ell_\alpha\rVert\leq2k_0/r$对一切$\alpha$成立。
\end{proof}

\paragraph{完备性用在哪里}
在$c_{00}$的例子中，$F_k$同样是覆盖空间的闭集，但没有一个含有开球：任何开球$B(x_0,r)$中都有$x_0+\frac r2\bm e_n$，取$n$足够大使它不属于$F_k$。论证失败的位置正是Baire纲定理，而Baire纲定理失败是因为空间不完备。同一个定理把线性、连续性和完备性组合起来，完成了量词的交换。

把$\ell_\alpha$换成取值在赋范空间中的有界线性算子，同样的四步论证给出算子族的一致有界原理；它与开映射定理一起，是第~\ref{chap:linear-operators}章研究完备空间中算子定理的出发点。本章只需要泛函的情形，下一节把它用于弱收敛列。

\section{弱收敛与弱星收敛}
\label{sec:lfunc-weak-convergence}

范数收敛要求整体误差$\lVert x_n-x\rVert$趋于零。在无限维空间中这个要求很强：第~\ref{chap:functions-as-vectors}章已经看到，单位球中的序列可以没有任何范数收敛的子列。换一个角度，只要求每一个固定的连续线性读数都收敛，会得到更弱、也更容易满足的收敛。本节分两种情形：在原空间$X$中用全部$\ell\in X^*$去读，得到弱收敛；在对偶空间$X^*$中只用原空间的向量去读，得到弱星收敛。两者的观测来源不同，不能只当作同一个概念的强弱两档。最后说明这两种收敛仍然控制序列的范数界。

\subsection{弱收敛}
\label{sec:lfunc-weak}

先用第~\ref{chap:topological-spaces}章的邻域语言描述“只看有限个读数时的接近”。给定$x_0\in X$、有限个泛函$\ell_1,\ldots,\ell_m\in X^*$和$\varepsilon>0$，令
\begin{equation}
U(x_0;\ell_1,\ldots,\ell_m;\varepsilon)=\{x\in X:\lvert\ell_i(x)-\ell_i(x_0)\rvert<\varepsilon,\ i=1,\ldots,m\}.
\label{eq:lfunc-weak-neighborhood}
\end{equation}
它由有限个读数的误差约束定义：只要这$m$个读数与$x_0$处的读数相差不到$\varepsilon$，就算“接近”。

\begin{definition}[弱拓扑与弱收敛]
\label{def:lfunc-weak-topology}
设$X$是赋范空间。若子集$G\subset X$对每个$x_0\in G$都包含某个形如式~\eqref{eq:lfunc-weak-neighborhood}的集合，称$G$为弱开集。弱开集组成的拓扑称为$X$上的\term{弱拓扑}（weak topology），记为$\sigma(X,X^*)$。序列$(x_n)$在弱拓扑中收敛到$x$，称为\term{弱收敛}（weak convergence），记为$x_n\rightharpoonup x$。
\end{definition}

需要核对弱开集确实组成拓扑。空集和$X$显然是弱开集，任意并的条件逐点检查即可。设$x_0$属于两个弱开集$G_1,G_2$之交，$G_1$含$U(x_0;\ell_1,\ldots,\ell_m;\varepsilon)$，$G_2$含$U(x_0;\ell'_1,\ldots,\ell'_{m'};\varepsilon')$；把两组泛函合在一起、误差取$\min(\varepsilon,\varepsilon')$，得到的集合同时含于两者，所以交仍是弱开集。式~\eqref{eq:lfunc-weak-neighborhood}本身也是弱开集：若$x_1$属于它，取$\delta=\varepsilon-\max_i\lvert\ell_i(x_1)-\ell_i(x_0)\rvert>0$，则$U(x_1;\ell_1,\ldots,\ell_m;\delta)$含于其中。

\begin{theorem}[弱收敛的读数判别]
\label{thm:lfunc-weak-criterion}
在赋范空间$X$中，$x_n\rightharpoonup x$当且仅当$\ell(x_n)\to\ell(x)$对每个$\ell\in X^*$成立。弱极限若存在则唯一。
\end{theorem}

\begin{proof}
设$x_n\rightharpoonup x$。对固定的$\ell$和$\varepsilon>0$，$U(x;\ell;\varepsilon)$是$x$的弱开邻域，序列最终落在其中，即$n$充分大时$\lvert\ell(x_n)-\ell(x)\rvert<\varepsilon$。反过来，设每个读数都收敛。任给$x$的弱开邻域，它包含某个$U(x;\ell_1,\ldots,\ell_m;\varepsilon)$；对每个$i$取$N_i$使$n\geq N_i$时$\lvert\ell_i(x_n)-\ell_i(x)\rvert<\varepsilon$，则$n\geq\max_iN_i$时$x_n$落在这个邻域中。这里只需对有限个$i$取最大值，所以不需要不同读数的收敛速度一致。

唯一性：若$x_n\rightharpoonup x$且$x_n\rightharpoonup y$，则$\ell(x)=\ell(y)$对一切$\ell\in X^*$成立，由定理~\ref{thm:lfunc-norming}得$x=y$。
\end{proof}

唯一性用到了连续线性读数足以区分向量，这是Hahn--Banach定理的结论；同样的理由说明弱拓扑是Hausdorff的。在Hilbert空间中，由Riesz表示定理，判别式变成$\langle y,x_n\rangle\to\langle y,x\rangle$对每个$y\in H$成立：弱收敛就是在每个固定方向上的内积收敛。

\begin{theorem}[范数收敛推出弱收敛]
\label{thm:lfunc-strong-implies-weak}
若$\lVert x_n-x\rVert\to0$，则$x_n\rightharpoonup x$。
\end{theorem}

\begin{proof}
对每个$\ell\in X^*$，$\lvert\ell(x_n)-\ell(x)\rvert\leq\lVert\ell\rVert\lVert x_n-x\rVert\to0$。
\end{proof}

\paragraph{逆向不成立}
在$\ell^2$中考察标准单位向量$\bm e_n$。两两距离为$\sqrt2$，没有范数收敛的子列。但对固定的$y\in\ell^2$，$\langle y,\bm e_n\rangle=\overline{y_n}$，由$\sum_n\lvert y_n\rvert^2<\infty$知它趋于零，所以$\bm e_n\rightharpoonup0$，而$\lVert\bm e_n\rVert=1$始终不变。更一般地，Hilbert空间中任何标准正交序列$(u_n)$都弱收敛到零，因为由Bessel不等式$\sum_n\lvert\langle y,u_n\rangle\rvert^2\leq\lVert y\rVert^2$，通项趋于零。

函数空间中的典型例子是振荡。$\{\sin(nx)/\sqrt\pi\}_{n\geq1}$是$L^2(0,2\pi)$中的标准正交序列，所以$\sin(nx)\rightharpoonup0$。取读数$g=\mathbf{1}_{[0,1]}$可以直接看到：$\int_0^1\sin(nx)\,dx=(1-\cos n)/n\to0$，而$\lVert\sin(nx)\rVert_2=\sqrt\pi$。函数在正负之间越摆越快，任何固定的权重与它相乘再积分，正负部分都相互抵消。弱收敛失去的正是对整体能量的控制：每个固定方向上的分量都在消失，能量却转移到越来越高的频率或越来越远的坐标上。

\paragraph{无限维中弱邻域很大}
弱邻域$U(0;\ell_1,\ldots,\ell_m;\varepsilon)$包含子空间$\ker\ell_1\cap\cdots\cap\ker\ell_m$，后者是$m$个线性方程的解集。若$X$是无限维的，映射$x\mapsto(\ell_1(x),\ldots,\ell_m(x))\in\mathbb F^m$不可能单射，这个交集含有非零向量，从而含有整条直线。所以无限维空间中每个弱邻域都是无界的，弱拓扑严格弱于范数拓扑。有限维时，坐标泛函有限个，弱邻域可以取成坐标误差都小于$\varepsilon$的集合，弱收敛与范数收敛一致。

\paragraph{部分读数不能代替全部读数}
弱收敛要求全部$\ell\in X^*$都收敛。只检查一部分读数，结论可能完全不同。在$\ell^2$中取$x_n=n\bm e_n$，每个坐标泛函$x\mapsto x_k$在$x_n$上的值最终为零，坐标逐个收敛到零向量。但取$y=(1/k)_k\in\ell^2$，$\langle y,x_n\rangle=1$不趋于零，所以$x_n$不弱收敛到零；由下面的定理~\ref{thm:lfunc-weak-bounded}，它无界，也不弱收敛到任何向量。坐标泛函的线性组合在$(\ell^2)^*$中稠密，可它们只是全部读数的一部分；缺少统一的范数界时，稠密子集上的收敛推不到全部读数上。

\subsection{弱星收敛}
\label{sec:lfunc-weak-star}

当变化的对象本身是读数，即$X^*$中的元素时，还有另一种自然的测试方法：只用原空间$X$中的向量去读它们。对$\ell_0\in X^*$、有限个向量$x_1,\ldots,x_m\in X$和$\varepsilon>0$，令
\[
U^*(\ell_0;x_1,\ldots,x_m;\varepsilon)=\{\ell\in X^*:\lvert\ell(x_i)-\ell_0(x_i)\rvert<\varepsilon,\ i=1,\ldots,m\}.
\]

\begin{definition}[弱星拓扑与弱星收敛]
\label{def:lfunc-weak-star}
设$X$是赋范空间。用上述集合代替式~\eqref{eq:lfunc-weak-neighborhood}，按定义~\ref{def:lfunc-weak-topology}的方式得到$X^*$上的拓扑，称为\term{弱星拓扑}（weak-star topology），记为$\sigma(X^*,X)$。序列$(\ell_n)$在其中收敛到$\ell$，称为\term{弱星收敛}（weak-star convergence），记为$\ell_n\overset{*}{\rightharpoonup}\ell$。
\end{definition}

与定理~\ref{thm:lfunc-weak-criterion}的证明完全相同，$\ell_n\overset{*}{\rightharpoonup}\ell$当且仅当$\ell_n(x)\to\ell(x)$对每个$x\in X$成立，即逐个向量的求值收敛。弱星极限唯一：若$\ell(x)=\ell'(x)$对一切$x$成立，则$\ell=\ell'$，这里不需要延拓定理。

\paragraph{与对偶空间自身的弱收敛比较}
$X^*$本身是赋范空间，它也有定义~\ref{def:lfunc-weak-topology}意义下的弱拓扑$\sigma(X^*,X^{**})$，测试对象是$X^{**}$的全部元素。每个$x\in X$通过$\Phi_x(\ell)=\ell(x)$给出$X^{**}$的一个元素（第~\ref{sec:lfunc-reflexivity}小节将它记为$J(x)$），所以弱星收敛只用$X^{**}$中这一部分元素测试，要求更少：
\[
\ell_n\rightharpoonup\ell\ \text{（在}X^*\text{中）}\quad\Longrightarrow\quad\ell_n\overset{*}{\rightharpoonup}\ell .
\]
逆向一般不成立，因为$X^{**}$中可以有不来自$X$的元素。

\begin{example}[弱星收敛而不弱收敛]
\label{ex:lfunc-weak-star-not-weak}
取$X=\ell^1$。由定理~\ref{thm:lfunc-lp-dual}（计数测度，$p=1$），$X^*$可以等同于$\ell^\infty$，$y\in\ell^\infty$读作$x\mapsto\sum_ky_kx_k$。令$y^{(n)}\in\ell^\infty$在第$n$项及以后为$1$、此前为$0$。对每个$x\in\ell^1$，$\sum_ky^{(n)}_kx_k=\sum_{k\geq n}x_k\to0$，所以$y^{(n)}\overset{*}{\rightharpoonup}0$。

但$y^{(n)}$在$\ell^\infty$中不弱收敛。取第~\ref{sec:lfunc-hahn-banach}小节构造的$\Lambda\in(\ell^\infty)^*$，它在收敛数列上取极限，于是$\Lambda(y^{(n)})=1$对一切$n$成立，不趋于$\Lambda(0)=0$。若$y^{(n)}$弱收敛到某个$y$，由上面的蕴含它也弱星收敛到$y$，由弱星极限唯一得$y=0$，矛盾。$\Lambda$正是$X^{**}$中不来自$X$的元素。
\end{example}

\subsection{弱收敛列的范数控制}
\label{sec:lfunc-weak-norm-control}

弱收敛只给出逐个读数的信息，但借助一致有界原理，它仍然控制了序列的范数。关键的想法是把向量看成对偶空间上的读数：$x\in X$对应$X^*$上的泛函$\ell\mapsto\ell(x)$，由式~\eqref{eq:lfunc-norm-dual}，这个泛函的范数恰为$\lVert x\rVert$。

\begin{theorem}[弱收敛列有界且范数下半连续]
\label{thm:lfunc-weak-bounded}
设$X$是赋范空间，$x_n\rightharpoonup x$。则$\sup_n\lVert x_n\rVert<\infty$，并且
\begin{equation}
\lVert x\rVert\leq\liminf_{n\to\infty}\lVert x_n\rVert .
\label{eq:lfunc-weak-lsc}
\end{equation}
\end{theorem}

\begin{proof}
\emph{有界性。}令$\Phi_n(\ell)=\ell(x_n)$，$\ell\in X^*$。$\Phi_n$是$X^*$上的有界线性泛函，由式~\eqref{eq:lfunc-norm-dual}，$\lVert\Phi_n\rVert=\lVert x_n\rVert$。对每个固定的$\ell$，$\Phi_n(\ell)=\ell(x_n)$收敛，因而有界，所以$\{\Phi_n\}$在$X^*$上逐点有界。由定理~\ref{thm:lfunc-dual-complete}，$X^*$是Banach空间，由定理~\ref{thm:lfunc-uniform-boundedness}得$\sup_n\lVert x_n\rVert=\sup_n\lVert\Phi_n\rVert<\infty$。

\emph{下半连续。}任取$\lVert\ell\rVert\leq1$，则$\lvert\ell(x)\rvert=\lim_n\lvert\ell(x_n)\rvert\leq\liminf_n\lVert x_n\rVert$。对$\ell$取最大值，由式~\eqref{eq:lfunc-norm-dual}得式~\eqref{eq:lfunc-weak-lsc}。
\end{proof}

证明中用到完备性的空间是$X^*$而不是$X$，所以结论对不完备的$X$也成立。式~\eqref{eq:lfunc-weak-lsc}可以是严格不等式：$\bm e_n\rightharpoonup0$，而$\lVert0\rVert=0<1$。弱极限允许一部分能量在极限中消失，但不会凭空增加。

\paragraph{弱星收敛列}
对弱星收敛，逐点有界直接发生在$X$上。若$X$是Banach空间且$\ell_n\overset{*}{\rightharpoonup}\ell$，则对每个$x$，$(\ell_n(x))$收敛因而有界，由定理~\ref{thm:lfunc-uniform-boundedness}得$\sup_n\lVert\ell_n\rVert<\infty$；同样有$\lVert\ell\rVert\leq\liminf_n\lVert\ell_n\rVert$，因为$\lvert\ell(x)\rvert=\lim\lvert\ell_n(x)\rvert\leq\liminf_n\lVert\ell_n\rVert$对单位向量$x$成立。这里$X$的完备性不能省：在第~\ref{sec:lfunc-uniform-boundedness}节的$c_{00}$中，$\ell_n(x)=nx_n$对每个$x$最终为零，所以$\ell_n\overset{*}{\rightharpoonup}0$，但$\lVert\ell_n\rVert=n$无界。

\paragraph{有界只是必要条件}
定理~\ref{thm:lfunc-weak-bounded}给出了一个检验：无界序列不可能弱收敛。例如$\sqrt n\,\bm e_n$在$\ell^2$中不弱收敛，可以直接取$y_k=k^{-3/4}$验证，$y\in\ell^2$而$\langle y,\sqrt n\,\bm e_n\rangle=n^{1/4}\to\infty$。但有界不能反推收敛：$x_n=(-1)^n\bm e_1$有界，在读数$\langle\bm e_1,\cdot\rangle$下交替取$\pm1$，整列不弱收敛。有界序列是否至少有弱收敛的子列，是下一节的紧性问题，答案取决于空间。

\section{闭球的紧性}
\label{sec:lfunc-ball-compactness}

有限维空间中，有界序列总有收敛子列，因为闭球是紧的。无限维空间中闭球在范数下不紧，这是第~\ref{chap:functions-as-vectors}章的结论。换成上一节较弱的拓扑，紧性可以恢复，但要分清三件事：闭球在哪个空间里，用的是哪个拓扑，以及紧性能否转化为“有界序列有收敛子列”。本节先证明对偶空间的闭球在弱星拓扑下紧，再借助反射性把它转回原空间闭球的弱紧性，最后讨论收敛子列的提取。

记$B_X=\{x\in X:\lVert x\rVert\leq1\}$，$B_{X^*}=\{\ell\in X^*:\lVert\ell\rVert\leq1\}$，分别是原空间与对偶空间的闭单位球。

\subsection{对偶闭球的弱星紧性}
\label{sec:lfunc-banach-alaoglu}

弱星拓扑只要求逐个向量上的读数接近，测试条件很少，开集也就很少；开集越少，开覆盖越难找到无法取有限子覆盖的反例，紧性越容易成立。下面的定理说，对偶空间的闭单位球在这个拓扑下总是紧的。

\begin{theorem}[Banach--Alaoglu定理]
\label{thm:lfunc-banach-alaoglu}
设$X$是赋范空间。则$B_{X^*}$在弱星拓扑$\sigma(X^*,X)$下紧。
\end{theorem}

这一结论称为\term{Banach--Alaoglu定理}（Banach--Alaoglu theorem）。一般情形的证明把每个$\ell\in B_{X^*}$看成一族数$(\ell(x))_{x\in X}$，其中$\lvert\ell(x)\rvert\leq\lVert x\rVert$，从而把$B_{X^*}$放进无穷多个闭圆盘的乘积中；弱星拓扑恰好是乘积拓扑在其上的限制，线性条件$\ell(\alpha x+\beta y)=\alpha\ell(x)+\beta\ell(y)$又使$B_{X^*}$成为乘积中的闭集。剩下的一步需要“任意多个紧空间的乘积仍紧”的Tychonoff定理，它超出了第~\ref{chap:topological-spaces}章所讨论的有限乘积，这里不给出一般情形的证明。$X$可分时，第~\ref{sec:lfunc-subsequences}小节将用度量和对角线论证给出完整证明。

\paragraph{推广到任意闭球}
对$\ell_0\in X^*$和$r>0$，映射$T(\ell)=\ell_0+r\ell$把$B_{X^*}$映成闭球$\{\ell:\lVert\ell-\ell_0\rVert\leq r\}$。它把弱星邻域$U^*(\ell;x_1,\ldots,x_m;\varepsilon)$恰好映成$U^*(T\ell;x_1,\ldots,x_m;r\varepsilon)$，逆映射同理，所以$T$是弱星拓扑下的同胚。紧集在连续映射下的像是紧的，于是$X^*$中任意有限半径的闭球都是弱星紧的。

\paragraph{定理断言了什么}
定理不要求$X$完备，也不要求任何其他条件。它断言的是对偶空间$X^*$的闭球在$\sigma(X^*,X)$下紧，既不是原空间$X$的闭球在弱拓扑下紧，也不是$B_{X^*}$在范数下紧。后者一般不成立：取$X=\ell^1$，$X^*=\ell^\infty$，$\ell^\infty$中的单位向量$\bm e_n$两两距离为$1$，没有范数收敛的子列，$B_{\ell^\infty}$在范数下不紧。

\subsection{反射性}
\label{sec:lfunc-reflexivity}

要把对偶空间上的结论带回原空间，需要一座桥。每个$x\in X$都给出$X^*$上的一个读数：用$x$去读一个泛函$\ell$，得到$\ell(x)$。

\begin{definition}[自然嵌入与反射空间]
\label{def:lfunc-reflexive}
对赋范空间$X$，令$J:X\to X^{**}$，
\[
J(x)(\ell)=\ell(x)\qquad(x\in X,\ \ell\in X^*),
\]
称为\term{自然嵌入}（canonical embedding）。若$J$是满射，称$X$为\term{反射空间}（reflexive space）。
\end{definition}

$J(x)$对$\ell$线性，且$\lvert J(x)(\ell)\rvert\leq\lVert\ell\rVert\lVert x\rVert$，所以$J(x)\in X^{**}$；$J$本身对$x$线性。由范数的对偶表示~\eqref{eq:lfunc-norm-dual}，
\[
\lVert J(x)\rVert=\sup_{\lVert\ell\rVert\leq1}\lvert\ell(x)\rvert=\lVert x\rVert,
\]
所以$J$是等距的，特别是单射。反射空间中$J$是等距同构，$X$与完备的$X^{**}$（定理~\ref{thm:lfunc-dual-complete}）等距对应，所以反射空间自动是Banach空间。定义要求的是这个特定的$J$为满射；仅仅存在某个从$X$到$X^{**}$的等距同构并不够，这里不展开相应的反例。

\paragraph{Hilbert空间是反射的}
记$R:H\to H^*$，$R(h)=\langle h,\cdot\rangle$。由Riesz表示定理，$R$是保范数的双射，并且共轭线性。给定$\Phi\in H^{**}$，映射$h\mapsto\overline{\Phi(Rh)}$对$h$线性（两次共轭抵消）且有界，再用一次Riesz表示定理，存在$x\in H$使$\overline{\Phi(Rh)}=\langle x,h\rangle$对一切$h$成立。于是
\[
\Phi(Rh)=\overline{\langle x,h\rangle}=\langle h,x\rangle=(Rh)(x)=J(x)(Rh).
\]
$R$是满射，所以$\Phi=J(x)$。

\paragraph{\texorpdfstring{$L^p$}{Lp}空间（\texorpdfstring{$1<p<\infty$}{1<p<infinity}）是反射的}
设测度$\sigma$有限，$q$为共轭指数，则$1<q<\infty$。由定理~\ref{thm:lfunc-lp-dual}，$(L^p)^*$的每个元素都是某个$L_g$，$g\in L^q$。给定$\Phi\in(L^p)^{**}$，映射$g\mapsto\Phi(L_g)$是$L^q$上的有界线性泛函，再对$L^q$用定理~\ref{thm:lfunc-lp-dual}，存在$f\in L^p$使$\Phi(L_g)=\int gf\,d\mu=L_g(f)=J(f)(L_g)$对一切$g$成立，所以$\Phi=J(f)$。两次表示都需要指数严格小于$\infty$，这就是端点被排除的原因。

\paragraph{完备而不反射}
$\ell^1$是Banach空间，却不反射。把$(\ell^1)^*$等同于$\ell^\infty$，则$J(x)(y)=\sum_ky_kx_k$。第~\ref{sec:lfunc-hahn-banach}小节的$\Lambda\in(\ell^\infty)^*=(\ell^1)^{**}$若等于某个$J(x)$，则$x_k=J(x)(\bm e_k)=\Lambda(\bm e_k)=0$对一切$k$成立，$x=0$，与$\Lambda(1,1,\ldots)=1$矛盾。所以完备性只说明Cauchy列有极限，反射性则要求二次对偶中没有多余的元素，后者严格更强。第~\ref{sec:lfunc-subsequences}小节还会看到$L^1(0,1)$也不反射。

\subsection{反射空间中的弱紧性}
\label{sec:lfunc-reflexive-weak-compact}

设$X$反射。$J$是等距的满射，所以它把$B_X$恰好映成$B_{X^{**}}$：$\lVert x\rVert\leq1$当且仅当$\lVert J(x)\rVert\leq1$，而$B_{X^{**}}$中每个元素都是某个$J(x)$。

\paragraph{两种拓扑的对应}
$X$上的弱拓扑$\sigma(X,X^*)$用有限个$\ell_i\in X^*$测试，$X^{**}=(X^*)^*$上的弱星拓扑$\sigma(X^{**},X^*)$也用有限个$\ell_i\in X^*$测试。由$J(x)(\ell_i)=\ell_i(x)$，
\[
J\bigl(U(x_0;\ell_1,\ldots,\ell_m;\varepsilon)\bigr)=U^*\bigl(J(x_0);\ell_1,\ldots,\ell_m;\varepsilon\bigr),
\]
右端用到了$J$为满射，否则只能得到右端与$J(X)$的交。两种拓扑的基本邻域一一对应，所以$J$是从$(X,\sigma(X,X^*))$到$(X^{**},\sigma(X^{**},X^*))$的同胚。

\begin{theorem}[反射空间中闭球的弱紧性]
\label{thm:lfunc-reflexive-weak-compact}
设$X$是反射空间。则$B_X$在弱拓扑$\sigma(X,X^*)$下紧；$X$中任意有限半径的闭球也是弱紧的。
\end{theorem}

\begin{proof}
对赋范空间$X^*$应用定理~\ref{thm:lfunc-banach-alaoglu}，$B_{X^{**}}$在$\sigma(X^{**},X^*)$下紧。由上面的对应，$B_X=J^{-1}(B_{X^{**}})$是这个紧集在同胚$J^{-1}$下的像，所以在弱拓扑下紧。平移与缩放$x\mapsto x_0+rx$在弱拓扑下是同胚，理由与第~\ref{sec:lfunc-banach-alaoglu}小节相同，所以任意闭球弱紧。
\end{proof}

由此，Hilbert空间和$1<p<\infty$时的$L^p$空间中，闭球都是弱紧的。这一结论不保证范数紧性：$\ell^2$的闭单位球弱紧，但$\bm e_n$没有范数收敛的子列。弱紧性换来的是更弱的收敛方式，不是无限维困难的消失。

\subsection{收敛子列的提取}
\label{sec:lfunc-subsequences}

紧性说的是开覆盖有有限子覆盖，我们真正想用的却常常是“有界序列有收敛子列”。在度量空间中两者等价，这是第~\ref{chap:metric-spaces}章的结论；但第~\ref{chap:topological-spaces}章已经说明，一般拓扑空间中紧性不能直接推出序列紧性。弱拓扑和弱星拓扑在无限维时一般不能由度量给出，所以上面的紧性定理还不能直接用来提取子列。

\paragraph{可分前对偶下弱星拓扑可度量}
设$X$可分，$\{w_k\}$是$X$的可数稠密子集。令$z_k=w_k/\max(1,\lVert w_k\rVert)$，则$\{z_k\}\subset B_X$，并且在$B_X$中稠密：若$\lVert x\rVert\leq1$且$\lVert x-w_k\rVert<\varepsilon$，则$\lVert w_k\rVert<1+\varepsilon$，从而$\lVert z_k-w_k\rVert<\varepsilon$，$\lVert x-z_k\rVert<2\varepsilon$。在$B_{X^*}$上令
\[
d(\ell,\ell')=\sum_{k=1}^{\infty}2^{-k}\lvert\ell(z_k)-\ell'(z_k)\rvert .
\]
每项不超过$2\cdot2^{-k}$，级数收敛。对称性与三角不等式逐项成立；若$d(\ell,\ell')=0$，则$\ell$与$\ell'$在$B_X$的稠密子集上相等，由连续性在$B_X$上相等，由线性在$X$上相等。所以$d$是$B_{X^*}$上的度量。

\begin{lemma}[球上弱星拓扑的度量]
\label{lem:lfunc-weak-star-metric}
设$X$可分。则在$B_{X^*}$上，度量$d$诱导的拓扑与弱星拓扑$\sigma(X^*,X)$限制到$B_{X^*}$上的拓扑相同。
\end{lemma}

\begin{proof}
只需说明每一点处，两种拓扑的邻域互相包含。

\emph{度量球包含弱星邻域。}给定$\ell_0\in B_{X^*}$和$r>0$，取$K$使$2^{1-K}<r/2$。若$\ell\in B_{X^*}$且$\lvert\ell(z_k)-\ell_0(z_k)\rvert<r/2$对$k\leq K$成立，则$d(\ell,\ell_0)<\sum_{k\leq K}2^{-k}r/2+\sum_{k>K}2\cdot2^{-k}<r/2+r/2=r$。

\emph{弱星邻域包含度量球。}给定$U^*(\ell_0;x_1,\ldots,x_m;\varepsilon)$，不妨设各$x_i\neq0$，令$c=\max_i\lVert x_i\rVert$，$u_i=x_i/c\in B_X$，只需使$\lvert\ell(u_i)-\ell_0(u_i)\rvert<\varepsilon/c$。对每个$i$取$k_i$使$\lVert u_i-z_{k_i}\rVert<\varepsilon/(4c)$。对$\ell\in B_{X^*}$，
\begin{equation*}
\begin{aligned}
\lvert\ell(u_i)-\ell_0(u_i)\rvert&\leq\lvert\ell(u_i-z_{k_i})\rvert+\lvert\ell(z_{k_i})-\ell_0(z_{k_i})\rvert+\lvert\ell_0(z_{k_i}-u_i)\rvert\\
&<\frac{\varepsilon}{2c}+\lvert\ell(z_{k_i})-\ell_0(z_{k_i})\rvert .
\end{aligned}
\end{equation*}
令$K=\max_ik_i$，$\delta=2^{-K}\varepsilon/(2c)$。若$d(\ell,\ell_0)<\delta$，则$\lvert\ell(z_{k_i})-\ell_0(z_{k_i})\rvert\leq2^{k_i}d(\ell,\ell_0)<\varepsilon/(2c)$，于是$\lvert\ell(u_i)-\ell_0(u_i)\rvert<\varepsilon/c$对每个$i$成立。
\end{proof}

统一的范数界在证明中起了关键作用：$\ell$与$\ell_0$都在单位球中，才能用有限个$z_k$上的接近控制任意$u_i$上的接近。在整个$X^*$上，这个度量一般不能给出弱星拓扑。

\begin{theorem}[可分前对偶下的弱星子列]
\label{thm:lfunc-weak-star-subsequence}
设$X$是可分赋范空间。则$X^*$中每个有界序列都有弱星收敛的子列。
\end{theorem}

\begin{proof}
设$\lVert\ell_n\rVert\leq C$，$\{z_k\}$如上。对每个$k$，$(\ell_n(z_k))_n$是有界数列。先抽子列使在$z_1$上的值收敛，再从中抽子列使在$z_2$上的值收敛，如此继续；取第$j$个子列的第$j$项组成对角子列$(\ell_{n_j})$，它从某项起是每个子列的子列，所以$\ell_{n_j}(z_k)$对每个$k$都收敛。

对$x\in B_X$及$\varepsilon>0$，取$z_k$使$\lVert x-z_k\rVert<\varepsilon$，则
\[
\lvert\ell_{n_i}(x)-\ell_{n_j}(x)\rvert\leq2C\varepsilon+\lvert\ell_{n_i}(z_k)-\ell_{n_j}(z_k)\rvert,
\]
最后一项在$i,j$充分大时小于$\varepsilon$。所以$(\ell_{n_j}(x))$是Cauchy数列，由线性对一切$x\in X$都收敛，极限记为$\ell(x)$。$\ell$是线性的，且$\lvert\ell(x)\rvert\leq C\lVert x\rVert$，故$\ell\in X^*$，$\ell_{n_j}\overset{*}{\rightharpoonup}\ell$。
\end{proof}

取$C=1$，极限仍在$B_{X^*}$中，所以$B_{X^*}$在度量$d$下序列紧；由第~\ref{chap:metric-spaces}章，序列紧的度量空间是紧的，再由引理~\ref{lem:lfunc-weak-star-metric}，这就是$B_{X^*}$的弱星紧性。可分时定理~\ref{thm:lfunc-banach-alaoglu}由此得到完整证明。

\paragraph{可分性不能省}
若前对偶不可分，序列结论可以失效，尽管Banach--Alaoglu定理仍保证紧性。$\ell^\infty$不可分；在$(\ell^\infty)^*$中考虑$\ell_n(x)=x_n$，$\lVert\ell_n\rVert=1$。对任意子列$(\ell_{n_j})$，取$x\in\ell^\infty$使$x_{n_j}=(-1)^j$、其余坐标为零，则$\ell_{n_j}(x)=(-1)^j$不收敛。所以这个有界序列没有弱星收敛的子列。紧与序列紧在这里确实分开了。

\paragraph{反射空间中的弱子列}
对弱拓扑有一个不需要可分性的结论。Eberlein--\v{S}mulian定理断言：在Banach空间中，一个子集的弱闭包弱紧，当且仅当该子集中的每个序列都有在$X$中弱收敛的子列。它的证明较长，这里不给出证明。与定理~\ref{thm:lfunc-reflexive-weak-compact}合起来得到：

\begin{theorem}[反射空间中的弱收敛子列]
\label{thm:lfunc-reflexive-subsequence}
设$X$是反射空间。则$X$中每个范数有界的序列都有弱收敛的子列。
\end{theorem}

\begin{proof}
有界序列含于某个闭球$B=\{x:\lVert x-x_0\rVert\leq r\}$。由式~\eqref{eq:lfunc-norm-dual}，$B=\bigcap_{\lVert\ell\rVert\leq1}\{x:\lvert\ell(x)-\ell(x_0)\rvert\leq r\}$，每个集合都是弱闭的，所以$B$弱闭，它的弱闭包就是自身；由定理~\ref{thm:lfunc-reflexive-weak-compact}，$B$弱紧。由Eberlein--\v{S}mulian定理，$B$中每个序列都有弱收敛的子列。
\end{proof}

于是在Hilbert空间和$1<p<\infty$的$L^p$空间中，有界序列总有弱收敛子列，不必另加可分性条件。$\bm e_n$本身就弱收敛到$0$；有界序列$(-1)^n\bm e_1$不收敛，但它的偶数项子列收敛。

\paragraph{非反射空间中可以失败}
在$L^1(0,1)$中令$f_n=n\mathbf{1}_{(0,1/n)}$，$\lVert f_n\rVert_1=1$。设某个子列$f_{n_j}\rightharpoonup f\in L^1$。对$0<a<1$，$g=\mathbf{1}_{(a,1)}\in L^\infty$给出$L^1$上的连续读数，而$n_j>1/a$时$\int_a^1f_{n_j}=0$，所以$\int_a^1f=0$。令$a\to0$，由控制收敛得$\int_0^1f=0$。但读数$g\equiv1$给出$\int_0^1f_{n_j}=1$，又得$\int_0^1f=1$，矛盾。所以这个有界序列没有弱收敛的子列。

质量始终为$1$，却向左端点越挤越窄；远离原点的读数看到零，整体积分看到一，没有$L^1$函数能同时满足。由定理~\ref{thm:lfunc-reflexive-subsequence}的逆否命题，$L^1(0,1)$不是反射空间。若把$f_n\,dt$看成$C([0,1])$上的读数$u\mapsto\int_0^1uf_n$，它们在逐个连续函数上收敛到$u(0)$：在更大的对偶空间中，极限以点质量的形式出现了。

\section{本章小结}
\label{sec:lfunc-summary}

本章开头问：哪些标量读数能由函数的范数稳定控制，它们长什么样，能否从局部扩展到全空间。泛函是从线性空间到标量的映射；加权积分在等价类上有定义，点评估只在由真正的函数组成的空间上有定义，平方范数则是非线性的。线性泛函连续当且仅当存在常数$C$使$\lvert\ell(x)\rvert\leq C\lVert x\rVert$，最小的$C$就是泛函范数。同一条读数规则是否连续、范数多大，取决于所配的范数：点评估在一致范数下范数为$1$，在积分范数下被窄尖峰破坏而不连续。

全部连续线性泛函组成对偶空间，无论原空间是否完备，对偶空间都是Banach空间，它的对偶是二次对偶。Hilbert空间中，每个连续线性泛函都是与唯一向量的内积，范数相等，证明只用正交投影而不要求可分；表示的前提是泛函已经连续、空间已经完备。$\sigma$有限测度下$1\leq p<\infty$时，$L^p$上的连续线性泛函都是以$L^q$函数为权重的积分，权重在几乎处处意义下唯一、范数相等；$p=\infty$时对偶更大。

一般空间中写不出表示，Hahn--Banach定理转而保证子空间上的有界线性泛函能保持范数地延拓到全空间，延拓一般不唯一。由此每个向量的范数等于单位泛函读数的最大值，不同向量能被连续线性读数区分，闭子空间外的点能被一个在子空间上为零的读数分开；$\ell^\infty$上还存在不来自$\ell^1$的连续线性泛函。完备性给出另一种统一控制：Banach空间上逐点有界的泛函族范数一致有界，不完备时可以失败。

弱收敛要求每个连续线性读数收敛，弱星收敛只要求在原空间的每个向量上求值收敛，二者的观测来源不同。范数收敛推出弱收敛，逆向在无限维中不成立，标准正交序列与快速振荡都弱收敛到零而范数不变；只检查一部分读数不能代替全部读数。对偶空间中的弱收敛推出弱星收敛，逆向需要二次对偶中没有多余元素。弱收敛列有界、范数下半连续，Banach前对偶上的弱星收敛列也有界；有界只是必要条件。

紧性方面，任意赋范空间的对偶闭球在弱星拓扑下紧；反射空间通过等距同构$J$把这一结论转回原空间，闭球弱紧。紧性要转成子列，还需要条件：前对偶可分时，对偶闭球上的弱星拓扑可度量，有界序列有弱星收敛子列；反射空间中由Eberlein--\v{S}mulian定理，有界序列有弱收敛子列，不需可分性。$L^1$中质量集中的序列说明非反射空间里这可以失败。这些都不是范数紧性。

本章的读数都输出一个标量。第~\ref{chap:linear-operators}章把输出推广为向量，研究有界线性算子、算子范数、伴随与完备空间中的算子定理；本章的对偶空间、Riesz表示和弱收敛，是那里定义伴随与讨论紧算子的工具。

